\documentclass[11pt,a4paper]{article}
\usepackage[margin=2.2cm]{geometry}
\usepackage{amsmath,amssymb}
\usepackage{graphicx}
\usepackage{booktabs}
\usepackage{xcolor}
\usepackage[colorlinks=true,linkcolor=blue!50!black,urlcolor=blue!50!black]{hyperref}
\usepackage{enumitem}
\usepackage{fontspec}
\setmonofont{Menlo}[Scale=0.82]
\setlist{itemsep=2pt,topsep=4pt}
\setlength{\emergencystretch}{3em}

\newcommand{\code}[1]{\texttt{#1}}
\newcommand{\jump}[1]{[\![#1]\!]}

\title{The extended ghost fluid method (xGFM) for Stokes flow in StagFDTools\\[2pt]
\large Viscosity jumps, compressibility and free surfaces on an unfitted staggered grid,
solved by GCR with a Cholesky preconditioner}
\author{StagFDTools, branch \texttt{iim-discontinuous-viscosity}}
\date{\today}

\begin{document}
\maketitle

\noindent This note documents an alternative to the immersed interface method (IIM, see
\code{docs/IIM\_notes}): an extended ghost fluid method (xGFM) in the spirit of
Liu, Fedkiw \& Kang (2000, J.~Comput.\ Phys.\ 160) and Egan \& Gibou (2020, J.~Comput.\ Phys.\
405). The question was whether a ghost-fluid type method gives $L^\infty$ convergence of the
pressure across a viscosity jump on an unfitted Cartesian grid (first order is enough) with
operators that are sparser and closer to symmetric than the IIM ones.
\begin{itemize}
  \item Section~\ref{sec:poisson}: a 2D Poisson prototype comparing GFM, xGFM and IIM.
  \item Sections~\ref{sec:jumps}--\ref{sec:solver}: the Stokes xGFM (jump structure, local
    closures, solver).
  \item Sections~\ref{sec:schmid}--\ref{sec:fs}: tests (incompressible and compressible
    inclusions, free surface).
  \item Section~\ref{sec:discussion}: relation to the IIM and to the StagFDTools stress form,
    open points.
\end{itemize}
The code is in \code{examples/\_Stokes\_IIM/xGFM\_Stokes.jl} (discretisation and solver) and
in the test scripts named in each section. Everything is first order; the aim is a robust
first-order sharp-interface method.

%---------------------------------------------------------------
\section{Poisson prototype}\label{sec:poisson}

\code{examples/\_Poisson\_Interface/} solves $-\nabla\cdot(\beta\nabla u)=f$ with piecewise
constant $\beta$ and jumps $\jump{u}=a$, $\jump{\beta\partial_n u}=b$ on a 5-point grid.
\begin{itemize}
  \item \textbf{GFM:} a face crossed by $\Gamma$ (fraction $\theta$ on the side of node $P$)
    gets the series coefficient $\hat\beta=1/(\theta/\beta_P+(1-\theta)/\beta_K)$, from a 1D
    model with continuous normal flux. The operator stays symmetric and 5-point, but the
    tangential part of the flux jump is ignored: on an oblique interface the face flux is
    $O(1)$ wrong.
  \item \textbf{xGFM:} the 1D model is corrected by the tangential flux jump
    $\jump{\beta}\,\partial_t u\,(\mathbf t\cdot\mathbf e)$, with $\partial_t u$ (continuous)
    traced from one-sided fits. The correction is lagged (fixed point; the operator stays
    the symmetric GFM one).
  \item \textbf{IIM:} orders 1 and 2, in the $u$-form and in the $z$-form ($z=\beta u$, as for
    Stokes).
\end{itemize}
Errors are measured per side and normalised by the field range on that side, otherwise the
stiff-side fluxes hide the soft-side errors.

\begin{figure}[h]
  \centering
  \includegraphics[width=\textwidth]{figures/Poisson_convergence.png}
  \caption{Poisson prototype (\code{Poisson\_Interface\_Convergence.jl}). $L^\infty$ errors of
  $u$ (top) and of the flux (bottom), worse side, for a circle in a uniform field (T1) and a
  flower-shaped interface with jumps $a,b\neq0$ (T2), contrast $10^3$ either way.}
  \label{fig:poisson}
\end{figure}

\noindent Results (Fig.~\ref{fig:poisson}):
\begin{itemize}
  \item GFM: flux error $O(1)$ (0.9--2.3 at all resolutions).
  \item xGFM: flux error $\approx O(h^{0.7\text{--}1})$ (soft inclusion T1: $9.2\times10^{-2}$ at
    $N=40$, $1.4\times10^{-2}$ at $N=640$) with the symmetric 5-point operator, after 25--55
    fixed-point iterations. $\partial_t u$ must be weighted towards the stiff side; the
    soft-side weighting diverges. For the stiff flower (T2) it is erratic up to $N=320$.
  \item IIM2 is the most accurate (flux $10^{-4}$--$10^{-5}$ at $N=640$) with 28--29 entries per
    row and an asymmetric operator. The $u$-form fails for a stiff inclusion, the $z$-form for
    a soft one.
\end{itemize}
So the xGFM gives first-order flux convergence with a symmetric compact base operator. The
lagged fixed point is not robust, which motivates the Krylov solver of Section~\ref{sec:solver}.

%---------------------------------------------------------------
\section{Stokes: jump structure}\label{sec:jumps}

On each side $s$ of $\Gamma$, with constant viscosity $\eta_s$ and bulk viscosity $\xi_s$
($\xi=\infty$: incompressible),
\[
\boldsymbol\sigma=-p\,\mathbf I+2\eta\left(\dot{\boldsymbol\varepsilon}-\tfrac13(\nabla\cdot\mathbf v)\mathbf I\right),\qquad
p=-\xi\,\nabla\cdot\mathbf v .
\]
Across $\Gamma$ the velocity and the traction are continuous: $\jump{\mathbf v}=0$,
$\jump{\boldsymbol\sigma}\mathbf n=0$. In the frame $(\mathbf t,\mathbf n)$ of $\Gamma$:
\begin{itemize}
  \item $\partial_t v_t$ and $\partial_t v_n$ are continuous (derivatives along $\Gamma$ of a
    continuous field);
  \item $\partial_n v_t$ and $\partial_n v_n$ jump: $\jump{\nabla\mathbf v}=\mathbf a\otimes\mathbf n$;
  \item the normal and shear tractions $N=\mathbf n\cdot\boldsymbol\sigma\mathbf n$ and
    $S=\mathbf t\cdot\boldsymbol\sigma\mathbf n$ are continuous.
\end{itemize}
With $M=\xi+\tfrac43\eta$ and $\lambda=\xi-\tfrac23\eta$, the constitutive law gives on each
side
\[
\partial_n v_n=\frac{N-\lambda\,\partial_t v_t}{M},\qquad
\partial_n v_t+\partial_t v_n=\frac{S}{\eta},
\]
and the full stress tensor
\begin{equation}\label{eq:sides}
\boldsymbol\sigma_s=N\,\mathbf n\otimes\mathbf n+S\,(\mathbf n\otimes\mathbf t+\mathbf t\otimes\mathbf n)
+\left(a_sN+b_s\,\partial_t v_t\right)\mathbf t\otimes\mathbf t,\qquad
a=\frac{\lambda}{M}=1-\frac{2\eta}{M},\quad b=M-\frac{\lambda^2}{M}=4\eta\left(1-\frac{\eta}{M}\right),
\end{equation}
\[
N=-\frac{M_s}{\xi_s}\,p_s-2\eta_s\,\partial_t v_t .
\]
In the incompressible limit $a=1$, $b=4\eta$, $M/\xi=1$, $1/M=0$:
\[
\boldsymbol\sigma_s=N\,\mathbf n\otimes\mathbf n+S\,(\mathbf n\otimes\mathbf t+\mathbf t\otimes\mathbf n)
+(N+4\eta_s\partial_t v_t)\,\mathbf t\otimes\mathbf t,\qquad p_s=-N-2\eta_s\partial_t v_t,
\]
so that $\jump{p}=-2\jump{\eta}\partial_t v_t$ and $\jump{\sigma_{tt}}=4\jump{\eta}\partial_t v_t$.

\paragraph{Consequence.} Once $N$, $S$ and $\partial_t\mathbf v$ are known, the stress on
\emph{either} side is known. $\partial_t\mathbf v$ is the only interface quantity that has to be
traced (one-sided linear fits, Landry's $\eta$-weighted average, as in the IIM). $N$ and $S$ are
the two ``fluxes'' that the ghost fluid method closes locally.

%---------------------------------------------------------------
\section{Discretisation}\label{sec:disc}

The unknowns are the standard ones: $V_x$, $V_y$ on the faces, $p$ at the centroids. The
stresses are evaluated once per centroid ($\sigma_{xx}$, $\sigma_{yy}$) and once per vertex
($\sigma_{xy}$), as in the standard staggered scheme, and the momentum rows are their
divergence. Where a stress location is not cut, the standard stencil is used. At cut locations
the closures below give, for the side $r$ of the momentum row that uses it, the stress
$\sigma^{(r)}$ as a linear form in the nodal unknowns. The two rows that share a stress
location see $\sigma^{(r)}$ of their own side; the two values differ by the jump of
Eq.~\eqref{eq:sides}, so that traction balance holds discretely.

\subsection*{Segment model}

Along a grid segment between two velocity nodes $a$, $b$ (a $V_x$ pair through a centroid or a
vertex, or a $V_y$ pair), $\mathbf v$ is taken piecewise linear with a kink where the tangent
line of $\Gamma$ (at the traced crossing) cuts the segment, at fraction $\ell$ from $a$. Each
Cartesian velocity derivative is written in the $(\mathbf t,\mathbf n)$ frame, e.g.
\[
\partial_x v_x=t_x^2\,\partial_t v_t+n_x^2\,\partial_n v_n+t_xn_x\,\frac{S}{\eta},\qquad
\partial_y v_x=t_xt_y\,\partial_t v_t+n_xn_y\,\partial_n v_n+t_xn_y\,\frac{S}{\eta}+(n_xt_y-t_xn_y)\,\partial_t v_n,
\]
and averaged over the segment with the side values:
\[
\frac{v_b-v_a}{h}=\ldots+t_xn_x\,\frac{S}{\hat\beta}+n_x^2\left(\frac{N}{\hat M}-\bar a\,\partial_tv_t\right),\qquad
\hat\beta=\Big(\frac{\ell}{\eta_a}+\frac{1-\ell}{\eta_b}\Big)^{-1},\quad
\frac1{\hat M}=\frac{\ell}{M_a}+\frac{1-\ell}{M_b},\quad \bar a=\ell a_a+(1-\ell)a_b .
\]
$\hat\beta$ is the series (harmonic) viscosity of the GFM; $\hat M$ and $\bar a$ only appear in
the compressible case.

\subsection*{Normal stresses at centroids}

$N$ is taken from the centroid pressure, $N=-(M_c/\xi_c)\,p_c-2\eta_c\,\partial_tv_t$ (side $c$
of the centroid). $S$ is eliminated with the segment relation of the $V_x$ pair (for
$\sigma_{xx}$; the $V_y$ pair for $\sigma_{yy}$). Only the product $2t_dn_dS$ appears, so there
is no division by $t_dn_d$:
\[
\sigma_{dd}^{(r)}=\left(n_d^2+t_d^2a_r-\frac{2\hat\beta n_d^2}{\hat M}\right)N
+2\hat\beta\,\frac{\Delta v_d}{h}
+\left(t_d^2b_r-2\hat\beta t_d^2+2\hat\beta n_d^2\bar a\right)\partial_tv_t .
\]
Incompressible: $\sigma_{dd}^{(r)}=-p_c+2\hat\beta\,\Delta v_d/h+\left(-2\eta_c-2\hat\beta(t_d^2-n_d^2)+4\eta_rt_d^2\right)\partial_tv_t$.
Without a crossing ($\eta$ uniform) this reduces to $-(\lambda/\xi)p+2\eta\,\partial_xv_x$, which
equals the standard deviatoric form $-p+2\eta(\partial_xv_x-\tfrac13\nabla\cdot\mathbf v)$ when
$\nabla\cdot\mathbf v=-p/\xi$ holds at that centroid. The standard form is therefore kept wherever
the four faces and the centroid are on one side.

\subsection*{Shear stress at vertices}

$\sigma_{xy}^{(r)}=n_xn_y\left(N(1-a_r)-b_r\,\partial_tv_t\right)+G\,S$ with
$G=n_xt_y+n_yt_x$. $S$ comes from one segment through the vertex: the $V_x$ pair ($y$-segment,
coefficient $t_xn_y$) if $n_y^2\ge n_x^2$, else the $V_y$ pair (coefficient $t_yn_x$). With this
choice $G/(t_xn_y)$ or $G/(t_yn_x)$ stays in $[0,1]$. In the incompressible case $N$ drops out
($1-a=0$, $1/M=0$). In the compressible case $N$ at the vertex is the mean of the $N$ of the
surrounding centroids (exact for a straight $\Gamma$ and piecewise-linear velocity).

\subsection*{Continuity at cut centroids}

The two segments through a centroid share $S$: $\hat\beta_x(\Delta v_x/h-A_x)=t_xn_xS$ and
$\hat\beta_y(\Delta v_y/h-A_y)=t_yn_yS=-t_xn_xS$, with
$A_d=t_d^2\partial_tv_t+n_d^2(N/\hat M_d-\bar a_d\partial_tv_t)$. Their sum eliminates $S$:
\[
w\left[\hat\beta_x\Big(\frac{\Delta v_x}{h}-A_x\Big)+\hat\beta_y\Big(\frac{\Delta v_y}{h}-A_y\Big)\right]=0,\qquad
w=\frac2{\hat\beta_x+\hat\beta_y}.
\]
Without a crossing this is $\nabla\cdot\mathbf v+p/\xi=0$.

\subsection*{Gravity}

With a body force $\rho_s\mathbf g$, $\partial_nN=-\rho\,\mathbf g\cdot\mathbf n$ differs on the two
sides. $N$ taken from a centroid at signed distance $d_c$ from $\Gamma$ is corrected as
$N^{(r)}=N_c+(\rho_c-\rho_r)\,(\mathbf g\cdot\mathbf n)\,d_c$ before it enters
$\sigma_{dd}^{(r)}$. This makes hydrostatics exact (flat interface between grid lines,
any position).

\subsection*{Properties}

\begin{itemize}
  \item \textbf{Exact for piecewise-linear solutions} (straight $\Gamma$, constant $p$ per side).
    Test (scratch scripts, $24\times20$ grid, $\Gamma$ at four orientations, LU):
    incompressible with $\eta_-/\eta_+=50,10^{-3},10^{3}$: $|\mathbf v-\mathbf v_{ex}|\le10^{-10}$,
    $|p-p_{ex}|\le5\times10^{-8}$; compressible with $(\eta,\xi)_-/(\eta,\xi)_+=(50,20)/(1,2)$,
    $(1,2)/(10^3,10^5)$, $(10^{-3},5\times10^{-4})/(1,2)$: $\le4.5\times10^{-11}$ (relative).
  \item \textbf{Width:} the $\partial_t\mathbf v$ forms come from fits of about 12 nodes, so cut rows
    have about 26 (incompressible) to 30 (compressible) entries, as many as IIM order~1.
  \item \textbf{Symmetric part:} without the $\partial_t\mathbf v$ terms and the compressible $N$ terms,
    the operator is the standard staggered Stokes operator with series viscosities
    $\hat\beta$ at cut centroids and vertices. Its velocity block is symmetric to
    $10^{-17}$; it is the preconditioner.
\end{itemize}

%---------------------------------------------------------------
\section{Solver}\label{sec:solver}

The full operator is non-symmetric (lagging the $\partial_t\mathbf v$ terms as in the Poisson
fixed point was not robust). It is solved by GCR (restarted, 30--100 directions), preconditioned
by one Powell--Hestenes step with the Cholesky factor of the symmetric part:
\[
\mathbf K_f=\mathbf K_{uu}^{PC}+\mathbf G\,\mathbf P_b\,\mathbf G^{T},\qquad
\mathbf P_b=(\mathbf C+\eta_b^{-1}\mathbf I)^{-1},\qquad
\mathbf s_u=\mathbf K_f^{-1}(\mathbf f_u-\mathbf G\mathbf P_b\mathbf f_p),\quad
\mathbf s_p=\mathbf P_b(\mathbf f_p+\mathbf G^T\mathbf s_u),
\]
with $\mathbf C=\mathrm{diag}(1/\xi)$ ($\mathbf P_b=\eta_b$ if incompressible, $\eta_b=10^3$).
This is the scheme used for the non-symmetric Jacobians of plasticity. Three details matter:
\begin{itemize}
  \item \textbf{Symmetric diagonal scaling} by $1/\sqrt{\eta}$ (velocities) and $\sqrt\eta$
    (pressures) of each node's side, so that all blocks are $O(1)$ whatever the contrast and the
    preconditioner stays symmetric. Without it GCR stops early at high contrast.
  \item \textbf{Lagrange multiplier} for the pressure level (incompressible, Dirichlet
    velocities). The weighted continuity rows of cut cells change the left null vector, so the
    singular system is slightly inconsistent and GCR stagnates near $10^{-9}$. With the
    multiplier ($\sum p=0$, or the mean over one side, e.g.\ the air) it converges to
    $10^{-12}$.
  \item \textbf{Compressible:} no null space, no multiplier; the deviatoric term of the
    preconditioner uses a common viscosity $\min(\hat\beta_x,\hat\beta_y)$ for $\sigma_{xx}$
    and $\sigma_{yy}$ of a cut centroid, which keeps $\mathbf K_{uu}^{PC}$ symmetric positive
    definite.
\end{itemize}
Iteration counts are nearly independent of $\eta_b$ ($10$--$10^5$). They are given with the
results below.

%---------------------------------------------------------------
\section{Incompressible inclusion (Schmid \& Podladchikov 2003)}\label{sec:schmid}

Circular inclusion of radius 2 in a box $[-5,5]^2$, pure shear, Dirichlet velocities from the
exact solution, contrast $10^{3}$ (stiff) and $10^{-3}$ (soft) (\code{xGFM\_Schmid2003.jl}); pressure fields are compared in Fig.~\ref{fig:fields} (\code{xGFM\_Schmid2003\_fields.jl}).
Pressure errors are taken after removing the mean.

\begin{figure}[h]
  \centering
  \includegraphics[width=\textwidth]{figures/xGFM_Schmid2003.png}
  \caption{Incompressible inclusion: $L^1$ (top) and $L^\infty$ (bottom) pressure errors.}
  \label{fig:schmid}
\end{figure}

\begin{center}
\begin{tabular}{llcccccc}
\toprule
case & method & $n=20$ & 40 & 80 & 160 & 320 & entries/row\\
\midrule
stiff & standard FD  & 1.33 & 1.72 & 7.06 & 9.54 & 8.38 & 11\\
      & xGFM         & 2.93 & 1.34 & 0.91 & 0.44 & 0.28 & 24--27\\
      & IIM stress 1 & 1.14 & 0.95 & 0.74 & 0.44 & 0.24$^*$ & 25--29\\
      & IIM stress 2 & 0.89 & 0.39 & 0.19 & 0.056 & 0.017$^*$ & 57--61\\
\midrule
soft  & standard FD  & 0.91 & 1.55 & 2.30 & 2.50 & 2.69 & 11\\
      & xGFM         & 1.04 & 0.70 & 1.28 & 0.46 & 0.27 & 24--27\\
      & IIM stress 1 & 0.93 & 0.68 & 0.50 & 0.26 & 0.16$^*$ & 25--29\\
      & IIM stress 2 & 0.25 & 0.19 & 0.066 & 0.042 & 0.0068$^*$ & 57--61\\
\bottomrule
\end{tabular}

\smallskip
$L^\infty(P)$. $^*$: solved by GCR with the xGFM preconditioner (Section~\ref{sec:iimpc}).
\end{center}

\noindent Discussion:
\begin{itemize}
  \item The xGFM converges in $L^\infty(P)$ where standard FD diverges, at about the level of
    IIM order~1. The $L^1$ error is the lowest of the first-order methods for the stiff
    inclusion ($n=160$: $3.9\times10^{-3}$ against $1.1\times10^{-2}$ for IIM stress~1).
  \item The error consists of isolated one-cell spikes at cut cells plus a smooth $O(h)$ error
    in the soft phase (Fig.~\ref{fig:errmaps}). The soft-inclusion bump at $n=80$ is one cell
    whose centroid lies almost exactly one cell from $\Gamma$.
  \item GCR iterations: 28--109, rising slowly with $n$, nearly independent of $\eta_b$.
\end{itemize}

\begin{figure}[h]
  \centering
  \includegraphics[width=\textwidth]{figures/xGFM_Schmid2003_errormaps.png}
  \caption{xGFM pressure error, soft (top) and stiff (bottom) inclusion, $n=40,80,160$.}
  \label{fig:errmaps}
\end{figure}

\begin{figure}[h]
  \centering
  \includegraphics[width=\textwidth]{figures/xGFM_Schmid2003_fields.png}
  \caption{Pressure at $n=80$, stiff (top) and soft (bottom) inclusion: analytical solution
  (Schmid \& Podladchikov, 2003) against standard FD (phase-averaged viscosity), IIM (stress form, order 2) and xGFM,
  on a common colour scale (black/magenta: above/below the analytical range); right, profile along the
  cell row just above $y=0$. FD has spurious extrema all around the interface; the xGFM error is
  confined to a few cut cells where the interface is tangent to the grid lines.}
  \label{fig:fields}
\end{figure}

\subsection*{The xGFM preconditioner for the IIM}\label{sec:iimpc}

The Cholesky factor of the symmetric xGFM part also preconditions GCR on the IIM stress-form
matrices: the same unknowns and the same scaling, the sign of the momentum rows flipped.
\begin{itemize}
  \item It reproduces the direct IIM solutions in 170--500 iterations.
  \item It makes $n=320$ affordable without 10~GB direct solves: IIM stress~2 gives
    $L^\infty(P)=1.7\times10^{-2}$ (stiff) and $6.8\times10^{-3}$ (soft) in 16--26~s.
  \item IIM stress~1 at $n=320$ (stiff) reached the 1000-iteration limit at a residual of
    $1.7\times10^{-8}$.
\end{itemize}
Second-order accuracy and the usual Cholesky-preconditioned Krylov strategy are therefore
compatible.

%---------------------------------------------------------------
\section{Compressible inclusion (Moutzouris 2026)}\label{sec:moutz}

Same setup as IIM step~6 (\code{xGFM\_Moutzouris2026.jl}): circle and ellipse (aspect ratio 2,
far field rotated by $\pi/6$), $\eta_m=1$, $\xi_m=2$, stiff inclusion
($\eta_i=10^3$, $\xi_i=10^5$) and weak inclusion ($\eta_i=10^{-3}$, $\xi_i=5\times10^{-4}$, so
$\lambda_i<0$), Dirichlet velocities from the exact solution. The pressure is determined, so
errors are absolute. GCR restart 100.

\begin{figure}[h]
  \centering
  \includegraphics[width=\textwidth]{figures/xGFM_Moutzouris2026.png}
  \caption{Compressible inclusion: $L^1$ (top) and $L^\infty$ (bottom) pressure errors.
  IIM up to $n=160$ (direct solves).}
  \label{fig:moutz}
\end{figure}

\begin{center}
\begin{tabular}{llccccc}
\toprule
case & method & $n=20$ & 40 & 80 & 160 & 320\\
\midrule
circle, stiff & standard FD & 1.29 & 1.65 & 4.44 & 6.01 & 5.38\\
              & xGFM        & 1.71 & 1.26 & 0.77 & 0.34 & 0.27\\
              & IIM stress 1& 0.98 & 0.79 & 0.51 & 0.25 & --\\
              & IIM stress 2& 0.56 & 0.27 & 0.11 & 0.037 & --\\
\midrule
circle, weak  & standard FD & 1.38 & 2.66 & 3.35 & 3.51 & 3.70\\
              & xGFM        & 0.83 & 0.90 & 0.52 & 0.37 & 0.23\\
              & IIM stress 1& 0.61 & 0.60 & 0.42 & 0.22 & --\\
              & IIM stress 2& 0.34 & 0.18 & 0.059 & 0.036 & --\\
\midrule
ellipse, stiff& standard FD & 6.26 & 6.08 & 6.57 & 6.55 & 8.04\\
              & xGFM        & 1.29 & 1.14 & 0.68 & 0.43 & 0.28\\
              & IIM stress 1& 0.58 & 0.76 & 0.48 & 0.40 & --\\
              & IIM stress 2& 0.83 & 0.44 & 0.24 & 0.090 & --\\
\midrule
ellipse, weak & standard FD & 2.64 & 3.01 & 3.12 & 3.43 & 3.67\\
              & xGFM        & 1.78 & 1.37 & 1.09 & 0.54 & 0.31\\
              & IIM stress 1& 1.51 & 1.04 & 0.66 & 0.40 & --\\
              & IIM stress 2& 0.55 & 0.28 & 0.099 & 0.050 & --\\
\bottomrule
\end{tabular}

\smallskip
$L^\infty(P)$, compressible inclusion.
\end{center}

\noindent Discussion:
\begin{itemize}
  \item Standard FD does not converge in $L^\infty$ (3--8). The xGFM converges at roughly
    $O(h^{0.7})$ in all four cases, with errors 1.3--1.7 times those of IIM order~1.
  \item $L^1(P)$ of the xGFM is the lowest of the first-order methods at $n\ge80$ in all four
    cases (e.g.\ circle stiff, $n=160$: $1.8\times10^{-3}$ against $5.5\times10^{-3}$ for IIM1 and
    $3.1\times10^{-2}$ for FD).
  \item GCR iterations (to $10^{-12}$): 40--130 for the stiff inclusion, 40--235 for the weak
    circle, but 260--900 for the weak ellipse. With restart 30 the weak ellipse did not converge
    in 500 iterations. The compressible preconditioner is less close to the operator: the cut
    centroids couple $p$ with the coefficient $(M/\xi)(n_d^2+t_d^2a-\ldots)$ (3.7 inside the weak
    inclusion) while the preconditioner keeps $-1$.
\end{itemize}

%---------------------------------------------------------------
\section{Free surface}\label{sec:fs}

Two ways of using the xGFM for a free surface were tested (\code{xGFM\_FreeSurface.jl}), on
the two free-surface tests of the IIM notes (steps~9 and~10).
\begin{itemize}
  \item \textbf{Sticky air:} the two-phase xGFM with a small $\eta_{air}$. The series viscosity
    of a cut segment goes to zero with $\eta_{air}$, so no harmonic smearing.
  \item \textbf{Void} (\code{AssembleXGFM(...; void=:pos)}): the limit $\eta_{air}=0$ of the
    closures, with no unknowns in the void.
    \begin{itemize}
      \item $\hat\beta=0$ on cut segments, so no void velocity enters a fluid row (checked in
        the assembly).
      \item A void stress point carries the traction-free stress, $N=S=0$ in
        Eq.~\eqref{eq:sides}, with the gravity correction $-\rho_f(\mathbf g\cdot\mathbf n)d$.
      \item Every cut fluid centroid gets the pressure condition $N_\Gamma=0$ extended to the
        centroid, $N_c+\rho(\mathbf g\cdot\mathbf n)d_c=0$, instead of the continuity row.
        Keeping the weighted continuity row where only one face pair is cut leaves the pressure
        of a column undetermined when $\Gamma$ runs just below a row of $V_y$ faces
        (hydrostatic test with the surface between a centroid and the face above: $p$ error 4.0
        with LU). This is the same failure as the IIM free surface without
        \code{faces=:adjacent}.
      \item No pressure null space, no Lagrange multiplier.
    \end{itemize}
    Hydrostatics is exact (round-off) for any position of a flat surface.
\end{itemize}

\subsection*{Inviscid inclusion (exact solution)}

The Schmid \& Podladchikov inviscid inclusion seen through $[-0.25,0.25]\times[-0.75,-0.25]$
(pure shear, Dirichlet sides, the top of the box is void or air), as in IIM step~10.

\begin{figure}[h]
  \centering
  \includegraphics[width=\textwidth]{figures/xGFM_FreeSurface_Schmid.png}
  \caption{Free surface, inviscid inclusion. Top: pressure ($L^1$, $L^\infty$) and velocity
  errors in the matrix. Bottom: pressure error maps at $n=160$. Sticky-air pressures are
  compared after removing the mean; void and IIM free-surface pressures are absolute.}
  \label{fig:fsschmid}
\end{figure}

\begin{center}
\begin{tabular}{lcccccc}
\toprule
method & $n=20$ & 40 & 80 & 160 & 320 & 640\\
\midrule
xGFM void                     & 0.56 & 0.34 & 0.19 & 0.12 & 0.060 & 0.029\\
xGFM sticky, $\eta_{air}=10^{-2}$ & 0.48 & 0.29 & 0.16 & 0.084 & 0.072 & \\
xGFM sticky, $\eta_{air}=10^{-6}$ & 0.36 & 0.52 & 0.31 & 0.19 & 0.081 & \\
FD sticky, $\eta_{air}=10^{-4}$   & 2.37 & 2.57 & 2.64 & 2.71 & 2.75 & \\
IIM free surface (cubic fits) & 0.11 & 0.015 & $2.6\times10^{-3}$ & $6.8\times10^{-4}$ & $1.7\times10^{-4}$ & \\
IIM sticky, $\eta_{air}=10^{-6}$  & 0.024 & $5.8\times10^{-3}$ & $7.5\times10^{-4}$ & $3.3\times10^{-4}$ & -- & \\
\bottomrule
\end{tabular}

\smallskip
$L^\infty(P)$ in the matrix, inviscid inclusion with a free surface.
\end{center}

\noindent Discussion:
\begin{itemize}
  \item \textbf{Void:} $L^\infty(P)$ converges at $O(h)$ (0.56 to 0.029 from $n=20$ to 640);
    velocity $L^\infty$ at $O(h)$ ($1.1\times10^{-2}$ to $2.2\times10^{-4}$ at $n=320$). 12--19 GCR
    iterations, independent of $n$. The error is concentrated in the cut cells (the pressure
    condition extends $N_\Gamma$ to the centroid without its normal derivative), with a smooth
    $10^{-3}$ background. The IIM free surface (cubic one-sided fits, ghost extension) is
    $O(h^2)$ and 70--350 times more accurate ($n=80$--320).
  \item \textbf{Sticky air does not converge uniformly in $\eta_{air}$.} The first-order
    traction errors of cut cells act on the air as an $O(1)$ force that is not a gradient. Air
    velocities grow like $1/\eta_{air}$ (190 at $\eta_{air}=10^{-6}$, $n=40$), and their
    back-traction on the fluid, $\eta_{air}v_{air}/h$, does not decrease with $h$. The error
    stalls at coarse grids ($\eta_{air}=10^{-6}$: $L^1(P)\approx0.05$ for $n\le80$) and
    convergence resumes only when $h$ is small enough; GCR iterations grow to 900. With
    $\eta_{air}=10^{-2}$ the method converges but saturates at the modelling error of the air
    ($L^1\approx3.5\times10^{-3}$). It is still far better than FD sticky air, which does not
    converge in $L^\infty$.
\end{itemize}

\subsection*{Gaussian topography under gravity}

The relaxation test of IIM step~9: $\eta=\rho=1$, $g=-1$, box $[-3,3]\times[-5,1.5]$, no slip
on all sides, surface $y=\exp(-x^2/0.25)$. Reference: the IIM free surface (cubic fits).

\begin{figure}[h]
  \centering
  \includegraphics[width=\textwidth]{figures/xGFM_FreeSurface_Gaussian.png}
  \caption{Gaussian topography, $n_x=192$: IIM free surface (top), xGFM void (middle) and their
  difference (bottom).}
  \label{fig:gauss}
\end{figure}

\begin{center}
\begin{tabular}{lcccc}
\toprule
& $n_x=48$ & 96 & 192 & 384\\
\midrule
$\min V_y$, IIM free surface & $-0.20025$ & $-0.20164$ & $-0.20186$ & $-0.20204$\\
$\min V_y$, xGFM void        & $-0.21045$ & $-0.20662$ & $-0.20435$ & $-0.20292$\\
$\max|V_x|$, IIM free surface & 0.05401 & 0.05442 & 0.05450 & 0.05459\\
$\max|V_x|$, xGFM void        & 0.05536 & 0.05535 & 0.05502 & 0.05478\\
GCR iterations, xGFM void     & 20 & 21 & 23 & 25\\
GCR iterations, xGFM sticky $10^{-6}$ & 37 & 145 & 404 & 1292\\
\bottomrule
\end{tabular}
\end{center}

\noindent The xGFM void solution matches the IIM free surface (Fig.~\ref{fig:gauss}). The
difference in the peak velocities decreases at $O(h)$: in $\min V_y$ it is $1.0\times10^{-2}$,
$5.0\times10^{-3}$, $2.5\times10^{-3}$ and $9\times10^{-4}$. The field differences are
concentrated in the surface cells. The sticky-air xGFM with $\eta_{air}=10^{-6}$ again has
very large air velocities (up to 290 at $n_x=48$) and poor GCR convergence. With
$\eta_{air}=10^{-3}$ it is close to the void solution at $n_x\ge192$.

%---------------------------------------------------------------
\section{Discussion}\label{sec:discussion}

\subsection*{xGFM and IIM order 1}

Both are first order and need the same trace, $\partial_t\mathbf v$ from one-sided fits. Both
give rows of about 26--30 entries. The differences are:
\begin{itemize}
  \item \textbf{Flux form.} The xGFM evaluates one stress per centroid and per vertex, shared by
    the two rows that use it. The IIM corrects each stencil pointwise from a Taylor extension of
    the other side.
  \item \textbf{Local closure.} The xGFM closes the normal fluxes $N$ and $S$ with a 1D series
    model on one grid segment and needs no jump of $p$ or $\nabla p$. In the IIM these jumps
    follow from the side equations.
  \item \textbf{Symmetric part.} The xGFM operator without its trace terms is the standard
    staggered operator with series viscosities. It is symmetric positive definite in the
    velocity block, so it can be factorised once and used as a preconditioner. Its Cholesky
    factor also preconditions the IIM matrices (Section~\ref{sec:iimpc}).
\end{itemize}

\subsection*{Stress form}

The xGFM is already a stress form. It can be written as a kinematic reconstruction followed by
the unchanged constitutive kernel:
\begin{itemize}
  \item From $(N, S, \partial_t\mathbf v)$, reconstruct the one-sided velocity gradient of the
    row's side:
    \[
    \partial_n v_n=\frac{N-\lambda\partial_tv_t}{M},\qquad
    \partial_n v_t=\frac{S}{\eta}-\partial_t v_n .
    \]
  \item Evaluate $\boldsymbol\tau=\mathbf D\,\dot{\boldsymbol\varepsilon}$ with the side's
    material.
\end{itemize}
For isotropic linear viscosity this is identical to the formulas of Section~\ref{sec:disc}.
For anisotropy or power-law creep, the series closure becomes a small local problem at each
crossing: find the one-sided strain rates for which the tractions computed with both sides'
constitutive laws are continuous. For anisotropy it couples $N$ and $S$ (a $2\times2$ system).
It is local, so it can be differentiated by AD for a Newton Jacobian. One caveat: a general
$\mathbf D$ needs $\dot\varepsilon_{xy}$ at centroids and $\dot\varepsilon_{xx}$,
$\dot\varepsilon_{yy}$ at vertices, so $S$ must be taken from the segment closer to the normal
at centroids as well. This is not implemented.

\subsection*{Open points}
\begin{itemize}
  \item \textbf{Accuracy.} Everything is first order. IIM order~2 is 5--40 times more accurate
    on the inclusions. A second-order xGFM would need the normal derivatives of $N$ and $S$
    (curvature, $\jump{\nabla p}$) in the 1D closures.
  \item \textbf{Free surface.} The void xGFM is $O(h)$ and robust: about 20 iterations,
    exact hydrostatics. The pressure condition in cut cells extends $N_\Gamma$ without its normal
    derivative; a linear extension (as the IIM \code{pghost=:linear}) should reduce the
    cut-cell error.
  \item \textbf{Sticky air with small $\eta_{air}$} needs consistent tractions on the air side.
    The void formulation is preferable.
  \item \textbf{Compressible preconditioner} (weak ellipse: up to 900 iterations).
  \item The code is in \code{examples/}. Moving it to \code{src/StokesIIM} would share the
    trace and numbering machinery with \code{AssembleIIMStress}.
\end{itemize}

\end{document}
